\section{Spin correlations and the full transfer spectrum}
\label{sec:sector}

A partition function counts every transfer eigenvalue, whereas a spin
two-point function sees only states created by a spin component.
We first relate these two measurements. The key estimate is
finite-dimensional: the covariance of two even spin monomials is
quadratic in the largest component two-point function between their
supports. The proof adapts the three-component argument of
\cite[Section~4]{O3Pole}, replacing its conditional Ising--rotator
decomposition by two conditional rotators. We include the proof because
its treatment of several spin components is essential for the full gap.

Here a monomial is a product of individual spin components with
coefficient one. Its site list has one entry for each factor,
including repetitions.

\subsection{Conditional rotators and correlation signs}

Let $\mathcal G=(V,E)$ be a finite graph and give $q_x\in S^3$, $x\in V$,
the probability law
\begin{equation}
 d\nu_J(q)=Z_J^{-1}
 \exp\left(\sum_{xy\in E}J_{xy}q_x\cdot q_y\right)
 \prod_{x\in V}d\omega_3(q_x),
 \qquad J_{xy}\ge0.
 \label{eq:wilson-graph-law}
\end{equation}
There is no external field and no pinned boundary. We write
$\mathbb E_J$ for this expectation.

We recall the finite-system inequalities used below. For a zero-field
ferromagnetic plane rotator with angles $\theta_x$, every two integer
vectors $k,l$ indexed by $V$ satisfy
\begin{equation}
 \mathbb E\cos(k\cdot\theta)\ge0,\qquad
 \operatorname{Cov}\bigl(\cos(k\cdot\theta),
                         \cos(l\cdot\theta)\bigr)\ge0.
 \label{eq:wilson-xy-ginibre}
\end{equation}
These are the plane-rotator inequalities of Ginibre
\cite{Ginibre1970}. For clarity, the needed version has a short
replica proof. Fourier expansion with nonnegative coefficients gives
the first assertion. For the second, take independent copies
$\theta,\theta'$ and substitute
$\theta=\varphi+\psi$, $\theta'=\varphi-\psi$. This surjective
homomorphism of product tori preserves Haar integration. Each
observable difference becomes
$-2\sin(k\cdot\varphi)\sin(k\cdot\psi)$, and the replicated weight is
\[
 \exp\left(2\sum_{xy\in E}J_{xy}
       \cos(\varphi_x-\varphi_y)\cos(\psi_x-\psi_y)\right).
\]
Expanding it expresses twice the covariance as a sum of squares of
real integrals with nonnegative coefficients. Absolute convergence
holds on the finite graph. Differentiating a cosine mean in a bond
coupling consequently shows that the mean is nondecreasing in that
coupling. We also use the zero-field Ising GKS inequalities
\cite{KellySherman1968}:
sign-monomial means and their covariances are nonnegative.
These are precisely the finite-system ingredients of the conditional
argument in \cite[Section~3.11]{SharpO4}.

A law of real variables is \emph{associated} if any two bounded
coordinatewise increasing functions have nonnegative covariance.
At every vertex write
\begin{equation}
 q_x=\bigl(\cos\alpha_x\cos\theta_x,\,
           \cos\alpha_x\sin\theta_x,\,
           \sin\alpha_x\cos\phi_x,\,
           \sin\alpha_x\sin\phi_x\bigr),
 \qquad 0\le\alpha_x\le\frac{\pi}{2}.
 \label{eq:wilson-two-planes}
\end{equation}
The sphere prior is the product of the two uniform circle laws and
$\sin(2\alpha_x)\,d\alpha_x$. Given $\alpha$, the two circle systems
are independent and have couplings
\[
 J^L_{xy}=J_{xy}\cos\alpha_x\cos\alpha_y,\qquad
 J^R_{xy}=J_{xy}\sin\alpha_x\sin\alpha_y.
\]

\begin{lemma}[Association and single-component signs]
\label{lem:wilson-signs}
The $\alpha$ marginal in \eqref{eq:wilson-two-planes} is associated.
For two even-degree monomials each using a single spin component,
the covariance is nonnegative when their components agree and
nonpositive when they differ. The same signs hold for the two
components of a zero-field ferromagnetic plane rotator.
\end{lemma}

\begin{proof}
The density of $\alpha$ relative to its product one-site prior is
proportional to $Z_L(\alpha)Z_R(\alpha)$. For either factor, denote
its bond observables by $O_e$ and its couplings by $J_e(\alpha)$.
At two distinct vertices,
\[
 \partial_x\partial_y\log Z
 =\sum_e\mathbb E(O_e)\,\partial_x\partial_yJ_e
 +\sum_{e,f}\operatorname{Cov}(O_e,O_f)
                  \partial_xJ_e\,\partial_yJ_f.
\]
All first coupling derivatives in the left plane are nonpositive;
all those in the right plane are nonnegative. A nonzero mixed
derivative at distinct vertices occurs only at a bond's endpoints
and is nonnegative. The bond means and covariances are nonnegative
by \eqref{eq:wilson-xy-ginibre}. Thus the logarithmic density has
nonnegative mixed derivatives. It satisfies the FKG lattice
condition and hence is associated \cite{FortuinKasteleynGinibre1971}.
For this continuous array, restrict to compact subintervals of
$(0,\pi/2)$, apply the finite-grid FKG inequality to approximating
densities, and pass to the limit. Bounded increasing functions follow
by approximation and the usual monotone-class extension. The endpoint
restrictions can then be removed.

For the plane-rotator assertion, write its components as
$(\varepsilon\cos\gamma,\varepsilon'\sin\gamma)$,
$0\le\gamma\le\pi/2$. Given $\gamma$, the signs are independent
ferromagnetic Ising systems. The same mixed-derivative calculation
with GKS proves association of $\gamma$. Including amplitude factors,
the first-component monomial mean is nonnegative and decreasing in
$\gamma$, while the second-component mean is nonnegative and
increasing. Conditional independence gives the nonpositive covariance
for unlike components. For like components, the conditional covariance
is nonnegative by GKS and the covariance of conditional means is
nonnegative by association.

In \eqref{eq:wilson-two-planes}, a component-one monomial has
conditional mean equal to its cosine amplitudes times a nonnegative
linear combination of cosine-character means. It is nonnegative and
decreasing in $\alpha$. A component-four monomial has a nonnegative
increasing conditional mean, by a fixed rotation of the right circle.
Conditional independence gives the nonpositive covariance between
these components. For like components, the conditional covariance
is nonnegative by \eqref{eq:wilson-xy-ginibre}; the conditional means
have the same monotonicity. Component symmetry supplies the remaining
choices of components.
\end{proof}

\subsection{The covariance comparison}

\begin{proposition}[Finite-graph sector comparison]
\label{prop:wilson-covariance}
For the law \eqref{eq:wilson-graph-law}, let $P,Q$ be spin monomials
of positive degrees $r,s$ with site lists $\mathcal X,\mathcal Y$,
and put
\[
 G=\max_{x\in\mathcal X,\ y\in\mathcal Y}\mathbb E_J(q_x^1q_y^1).
\]
Then $G\ge0$, and a constant $C_{r,s}$ depending only on the degrees
satisfies
\begin{equation}
 |\operatorname{Cov}_J(P,Q)|\le C_{r,s}
 \begin{cases}
  G^2,&r,s\ \text{even},\\
  G,&r,s\ \text{odd},\\
  0,&r+s\ \text{odd}.
 \end{cases}
 \label{eq:sector-cov}
\end{equation}
If either monomial is constant, its covariance is zero. The estimate
is independent of the finite graph and its nonnegative couplings.
\end{proposition}

\begin{proof}
We first control coefficients in the positive pairing expansion.
This treats the odd case and several even cases. For the remaining
case we must also control the covariance of conditional means.

The uniform sphere moment tensor vanishes in odd degree, and
in degree $2h$ it is
\[
 \int_{S^3}q^{i_1}\cdots q^{i_{2h}}\,d\omega_3(q)
 =\frac{1}{4\cdot6\cdots(4+2h-2)}
   \sum_{\pi}\prod_{\{b,c\}\in\pi}\mathbf1_{\{i_b=i_c\}},
\]
where the sum runs over pairings of labeled factors. Expand each
bond exponential in powers of the dot product and apply the tensor
identity at each site. Following the contracted bond lines pairs
the external factors, and every closed line contributes four.
Thus, for arbitrary assigned external components,
\begin{equation}
 \mathbb E_J\prod_{b=1}^{r+s}q_{x_b}^{i_b}
 =\sum_\pi d_\pi
       \prod_{\{b,c\}\in\pi}\mathbf1_{\{i_b=i_c\}},
 \qquad d_\pi\ge0.
 \label{eq:wilson-pairings}
\end{equation}
The $d_\pi$ do not depend on the component assignment.
Repeated sites cause no difficulty because the factors remain
labeled. Absolute convergence justifies regrouping the expansion.
In particular the component two-point functions are nonnegative.

A pair is \emph{crossing} if it joins $\mathcal X$ to $\mathcal Y$.
When $r,s$ are odd, every pairing crosses. Fix one crossing pair
and assign component four to its endpoints and component one to
all remaining endpoints. The resulting moment contains $d_\pi$
as a nonnegative summand. Conditional on $\alpha$, the remaining
left-plane monomial has mean in $[0,1]$, including all amplitudes.
Its removal bounds the moment by the chosen component-four pair
mean. Hence $d_\pi\le G$. Both separate monomial means vanish
in odd total degree, so summing the pairing coefficients proves
the odd case.

When $r,s$ are even, every crossing pairing has at least two
crossing pairs. Assign component four to one, component two to
another, and component one to the remaining endpoints.
In the left plane, Lemma~\ref{lem:wilson-signs} bounds the joint
mean of the component-two pair and the component-one monomial
by the product of their nonnegative means. The latter monomial
mean and its amplitudes are at most one. The surviving
component-two pair mean, including its amplitudes, is decreasing
in $\alpha$, because
\[
 \mathbb E(\sin\theta_x\sin\theta_y\mid\alpha)
 =\tfrac12\mathbb E(\cos(\theta_x-\theta_y)\mid\alpha).
\]
This also holds when $x=y$. The component-four pair mean is
increasing in $\alpha$. Association bounds their product
expectation by the product of the full pair means, each at most $G$.
Therefore
\begin{equation}
 d_\pi\le G^2
 \quad\text{for every crossing pairing of two even lists}.
 \label{eq:wilson-crossing}
\end{equation}

For two single-component even monomials, compare equal and
different choices of their components. The products of their
means agree by symmetry. Their joint moments differ by exactly
the crossing contributions in \eqref{eq:wilson-pairings}.
By Lemma~\ref{lem:wilson-signs}, the two covariances have opposite
weak signs, so each absolute covariance is at most this difference.
Equation~\eqref{eq:wilson-crossing} proves the bound. The same
reasoning applies to every pair of even sublists with the
original value of $G$. If $P$ or $Q$ has odd degree in some
component, its separate mean is zero and every compatible
joint pairing crosses. This also follows from
\eqref{eq:wilson-crossing}.

It remains to consider monomials even in every component.
We use single-component replacements, whose full covariances are already
controlled, to bound both terms in the conditional covariance decomposition.
Write $P=P_LP_R$ by its two planes and put
\[
 b_{P,L}=\mathbb E(P_L\mid\alpha),\qquad
 b_{P,R}=\mathbb E(P_R\mid\alpha).
\]
Let $P_L^*$ replace all its components by component one and
$P_R^*$ replace all its components by component three. Denote
their conditional means by $B_{P,L},B_{P,R}$ and use the
corresponding notation for $Q$. Then
\[
 |b_{P,L}|,|b_{P,R}|\le1,\qquad
 0\le B_{P,L},B_{P,R}\le1,
\]
with $B_{P,L}$ decreasing and $B_{P,R}$ increasing in $\alpha$.

Factor the common amplitude $A_P(\alpha)$ from the left plane.
Its Fourier expansion and the expansion after replacement are
\[
 P_L=A_P(\alpha)\sum_k c_k\cos(k\cdot\theta),\qquad
 P_L^*=A_P(\alpha)\sum_k d_k\cos(k\cdot\theta),\qquad
 |c_k|\le d_k.
\]
Each sine or cosine supplies two exponential terms of modulus
$1/2$. Since the total number of sine factors is even, the
original coefficients are real and invariant under $k\mapsto-k$.
Replacement changes their phases to positive coefficients.
Collecting equal frequencies preserves their domination,
including for repeated sites. The same argument applies in
the right plane. By \eqref{eq:wilson-xy-ginibre},
\begin{equation}
 |\operatorname{Cov}(P_L,Q_L\mid\alpha)|
 \le\operatorname{Cov}(P_L^*,Q_L^*\mid\alpha),
 \label{eq:wilson-fourier-covariance}
\end{equation}
and likewise for $R$.

Conditional independence gives
\begin{align*}
 \operatorname{Cov}(P,Q\mid\alpha)
 ={}&\operatorname{Cov}(P_L,Q_L\mid\alpha)
                    \mathbb E(P_RQ_R\mid\alpha)\\
 &+b_{P,L}b_{Q,L}\operatorname{Cov}(P_R,Q_R\mid\alpha).
\end{align*}
The other factors have modulus at most one, so its averaged
absolute value is at most
\[
 \mathbb E\operatorname{Cov}(P_L^*,Q_L^*\mid\alpha)
 +\mathbb E\operatorname{Cov}(P_R^*,Q_R^*\mid\alpha).
\]
Each term is at most its full single-component covariance:
the omitted covariance of conditional means is nonnegative
by association. Both terms have already been bounded by
$C_{r,s}G^2$.

We still need the covariance of the conditional means.
The Fourier comparison shows that $B_{P,L}\pm b_{P,L}$ are
nonnegative combinations of cosine-character means multiplied
by their common amplitude. Both are decreasing; likewise
$B_{P,R}\pm b_{P,R}$ are increasing. For
$\alpha\le\widetilde\alpha$ this gives
\[
 |b_{P,L}(\widetilde\alpha)-b_{P,L}(\alpha)|
 \le B_{P,L}(\alpha)-B_{P,L}(\widetilde\alpha),
\]
and the corresponding increasing bound for the right plane.
The factors $b_{P,R},b_{P,L}$ in the conditional mean need not be monotone.
Their changes are nevertheless
bounded by those of the increasing $B_{P,R}$ and decreasing $B_{P,L}$.
To combine these two controls, put
\[
 Y_P=b_{P,R}b_{P,L},\qquad D_P=B_{P,R}-B_{P,L}.
\]
Using the modulus bounds on the $b$ functions yields
\[
 |Y_P(\widetilde\alpha)-Y_P(\alpha)|
 \le D_P(\widetilde\alpha)-D_P(\alpha).
\]
Consequently $D_P+Y_P$ and $D_P-Y_P$ are increasing.
Apply association to these functions and to $D_Q\pm Y_Q$.
Adding the two matching-sign inequalities and the two
opposite-sign inequalities proves
\begin{equation}
 |\operatorname{Cov}(Y_P,Y_Q)|
 \le\operatorname{Cov}(D_P,D_Q).
 \label{eq:wilson-conditional-means}
\end{equation}
On expanding the right side, the two same-plane terms are
bounded by their full single-component covariances, as above.
The two opposite-plane terms are exactly full single-component
covariances by conditional independence; they enter with minus
signs. Each is bounded in absolute value by the estimate already
proved on even sublists. Empty sublists give constants and zero
covariances. The total-covariance identity now proves the even case.

Finally, global spin inversion makes the joint moment zero
when $r+s$ is odd, and at least one separate mean is zero.
This proves all cases of \eqref{eq:sector-cov}.
\end{proof}

\subsection{Consequences for the plane and finite cylinders}

An observable is \emph{even} or \emph{odd} according to its parity
under simultaneous inversion $q_x\mapsto-q_x$ at every site.
A \emph{component reflection} reverses one fixed component at
every site and leaves the others unchanged. The same definitions
apply on a row. Thus total inversion is the product of the four
commuting component reflections; being even under total inversion
does not require being even under each of them.

For the original Wilson law, we write $q=\sigma$ and
$m(\beta)=m_{\mathrm{lat}}(\beta)$. Let $r_*(\beta)$ be the
supremum of the support of the $T_\beta$-spectral measure of
$[q_0^1]$ in the OS Hilbert space of Section~\ref{sec:introduction}.

\begin{proposition}[The spin field sees the full gap]
\label{prop:wilson-full-gap}
For every $\beta>0$, in the periodic plane state,
\begin{equation}
 m(\beta)=-\log r_*(\beta)
 =-\log\lim_{n\to\infty}
       \bigl(\mathbb E_{\mu_\beta}q_0^1q_{(n,0)}^1\bigr)^{1/n}.
 \label{eq:full-gap}
\end{equation}
The transfer norm on the even subspace orthogonal to the vacuum
is at most $r_*(\beta)^2$. Its gap is therefore at least
$2m(\beta)$. The logarithms can be read in the extended sense;
the imported upper bound below makes them finite throughout
the large-$\beta$ regime.
\end{proposition}

\begin{proof}
The finite-graph estimate passes to the periodic plane limit
for fixed monomials. Spatial translations are unitary in the
OS form and commute with $T_\beta$. Translated one-spin vectors
have the same spectral support. Thus OS Cauchy--Schwarz gives
\[
 0\le\mathbb E_{\mu_\beta}
       (q_{(0,x)}^1q_{(n,y)}^1)\le\tfrac14r_*^n
\]
for arbitrary spatial sites $x,y$. The factor $1/4$ is their
one-site reflection norm squared.

A reflected polynomial transfer autocorrelation is a finite
sum of monomial covariances at temporal separation at least
$n$. For an odd polynomial, \eqref{eq:sector-cov} bounds it by
$C_P r_*^n$; for a centered even polynomial it bounds it by
$C_P r_*^{2n}$. The spectral measures are positive, so their
supports lie in $[0,r_*]$ and $[0,r_*^2]$, respectively.
Polynomials on a fixed finite spin product are dense in the
$L^2$ space of its marginal law. Reflection invariance gives
$\|[F]\|_{\mathrm{OS}}^2\le\mathbb E|F|^2$, so this approximation
also holds in the OS seminorm. Centering and parity
symmetrization preserve density. Bounded spectral projections
extend the spectral bounds to the full subspaces.
Since $[q_0^1]$ is centered, the full nonvacuum norm is exactly
$r_*$. Finally, the $n$th root of a moment of a positive
spectral measure tends to the supremum of its support.
This proves \eqref{eq:full-gap}.
\end{proof}

For even integers $n,\ell\ge4$, let $Z_\beta(n,\ell)$ denote the Wilson partition function on
the rectangular torus with $n$ time sites and $\ell$ spatial
sites, with the same single-counted bonds and probability spin
measure as in Section~\ref{sec:introduction}.
For an even spatial circumference $\ell\ge4$, let
$\mathscr X_\ell=(S^3)^\ell$ with product probability measure.
Put $V_\ell(s)=\sum_{i=1}^{\ell}s_i\cdot s_{i+1}$, periodically,
and define the row transfer kernel
\begin{equation}
 K_{\beta,\ell}(s,s')
 =\exp\left[\frac{\beta}{2}V_\ell(s)
       +\beta\sum_{i=1}^{\ell}s_i\cdot s'_i
       +\frac{\beta}{2}V_\ell(s')\right].
 \label{eq:wilson-row-kernel}
\end{equation}
It is compact, with a continuous strictly positive kernel
and a simple top eigenvalue $\lambda_1$.
The tensor-power expansion of $e^{\beta s\cdot s'}$
makes it positive semidefinite and trace class. Write
\[
 U_{\beta,\ell}=K_{\beta,\ell}/\lambda_1,\qquad
 \Omega_{\beta,\ell}>0,\quad\|\Omega_{\beta,\ell}\|_2=1,\qquad
 \tau_{\beta,\ell}=\lambda_2/\lambda_1.
\]
For $\beta>0$ the kernel is not rank one, as its cross ratios
vary with the row configurations. Hence $0<\tau_{\beta,\ell}<1$.
Closing time gives
$Z_\beta(n,\ell)=\operatorname{Tr}K_{\beta,\ell}^{\,n}$.
These are the conventions of \cite[Section~2.1]{SharpO4}.
Denote the infinite-time cylinder expectation by
$\omega_{\beta,\ell}$.

\begin{lemma}[An odd detector for the first cylinder excitation]
\label{lem:wilson-detector}
The restriction of $U_{\beta,\ell}$ to the centered even row
subspace has norm at most $\tau_{\beta,\ell}^2$. There is a real
row function $F$, with $|F|\le1$, odd under reflection of one
spin component, such that
\begin{equation}
 \omega_{\beta,\ell}(F\,F\circ\vartheta_n)
 \ge\frac{\tau_{\beta,\ell}}{1+\tau_{\beta,\ell}}\,
                    \tau_{\beta,\ell}^{\,n},\qquad n\ge0,
 \label{eq:wilson-detector-overlap}
\end{equation}
where $\vartheta_n$ translates by $n$ rows.
\end{lemma}

\begin{proof}
The finite-graph estimate passes to the infinite-time cylinder
by the row transfer representation. Component vectors
$q_x^i\Omega$ are centered, so their correlations across $n$
rows are at most $\tau^n$, suppressing $\beta,\ell$.
The even polynomial estimate therefore places the spectral
support of every centered even polynomial vector in $[0,\tau^2]$.
Polynomials are uniformly
dense in continuous row functions, and the continuous positive
$\Omega$ has a positive minimum. Thus polynomials times $\Omega$
are dense in row space, proving the even-sector norm bound.

Choose a real unit $\tau$-eigenfunction $v$. Since
$\tau>\tau^2$, its eigenspace is odd under simultaneous inversion.
The four component reflections commute with one another and
with $U$, so $v$ can have definite parity under each; at least
one is odd. Relabel it as component one. The modulus $|v|$
is even. With $b_0=\langle\Omega,|v|\rangle^2$, positivity of
the kernel and the even-sector bound give
\[
 \tau\le\langle|v|,U|v|\rangle
       \le b_0+(1-b_0)\tau^2,
 \qquad b_0\ge\frac{\tau}{1+\tau}.
\]
Take $F=\operatorname{sgn}v$, equal to zero on its zero set.
Its overlap with $v$ is
$\langle v,F\Omega\rangle=\langle|v|,\Omega\rangle$.
Keeping this contribution in the positive spectral expansion
of its correlation proves \eqref{eq:wilson-detector-overlap}.
\end{proof}

The detector can involve the entire row. We next recover a
component correlation with a cost polynomial in the circumference.

\begin{lemma}[Visibility through component correlations]
\label{lem:wilson-visibility}
For fixed $\beta>0$ there is a finite constant $C_\beta$,
independent of even $\ell\ge4$, such that bounded row functions
$F,G$ of modulus at most one, both odd under reflection of
component one, satisfy
\begin{equation}
 |\omega_{\beta,\ell}(F\,G\circ\vartheta_n)|
 \le C_\beta\sum_{x\in R_0,\ y\in R_n}
                      \omega_{\beta,\ell}(q_x^1q_y^1),
 \qquad n\ge1.
 \label{eq:wilson-visibility}
\end{equation}
Here $R_0,R_n$ are the two rows of $\ell$ sites.
\end{lemma}

\begin{proof}
On a finite torus, write
\[
 q_i=(\varepsilon_i u_i,\sqrt{1-u_i^2}\,l_i),\qquad
 \varepsilon_i\in\{-1,1\},\quad u_i\in[0,1],\quad l_i\in S^2.
\]
The prior is a product: $\varepsilon_i$ is fair, $l_i$ is
uniform, and $u_i$ has density proportional to $\sqrt{1-u_i^2}$.
Given $u,l$, the signs are a ferromagnetic Ising model with
coupling $\beta u_i u_j$. In its Edwards--Sokal representation
\cite{EdwardsSokal1988}, an edge between equal signs is open
with probability $1-e^{-2\beta u_i u_j}$, and an edge between
different signs is closed. Given the edges, cluster signs
are independent and fair. If no cluster
joins the detector rows, flip all clusters touching the first
row. This reverses $F$ and leaves $G$ unchanged. Its conditional
product mean is therefore zero. The absolute detector
correlation is bounded by the row-connection probability,
and then by the sum of pair-connection probabilities.

For a pair $x,y$ the joint spin--edge representation gives
\begin{equation}
 \mathbb E(q_x^1q_y^1)
 =\mathbb E[u_xu_y\mathbf1_{\{x\leftrightarrow y\}}].
 \label{eq:wilson-fk-pair}
\end{equation}
Condition on all edges, directions $l$, and amplitudes other
than $u_i$. After summing the signs, each closed incident edge
has factor $e^{-J}$ and each open edge has factor $e^J-e^{-J}$,
where $J=\beta u_i u_j$. The cluster multiplicity depends only
on the fixed edges. Raising $u_i\in[0,1/4]$ to
$u_i+1/4\in[1/4,1/2]$ increases every open factor. Closed-edge
factors and perpendicular factors
\[
 \exp\bigl(\beta\sqrt{1-u_i^2}\sqrt{1-u_j^2}\,l_i\cdot l_j\bigr)
\]
change by bounded ratios depending only on $\beta$ and degree
four. The amplitude prior is bounded above and below on these
intervals. Comparing their integrals proves
\[
 \mathbb P(u_i\ge1/4\mid\text{all other data})\ge c_\beta>0
\]
uniformly in the volume. Applying this twice by the tower
property gives
$\mathbb E[u_xu_y\mid\text{edges},l]\ge c_\beta^2/16$.
The connection event is edge-measurable, so
\eqref{eq:wilson-fk-pair} bounds its probability by
$16c_\beta^{-2}\mathbb E(q_x^1q_y^1)$. Summing proves the
finite-torus estimate.

At fixed circumference the transfer representation converges
to the cylinder for all bounded Borel row insertions: separate
the ground projection, use their bounded multiplication norms,
and let the excited trace tail tend to zero. The estimate
therefore passes to the cylinder, including the sign detector.
\end{proof}

\section{Cylinder masses and a partition-function rate}
\label{sec:cylinders}

We now express the plane mass as the decay rate of a
volume-canceling combination of square partition functions.
Trace estimates will be uniform over the admitted cutoff
trajectories. In contrast, the visibility argument will be
applied only after fixing one trajectory and sending the
circumference to infinity.

\subsection{The imported Wilson trajectory bounds}

Here is the precise input from \cite{SharpO4}. Fix its
sufficiently large dyadic blocking factor $L$ and auxiliary
parameters, and then choose the terminal reference coupling
$H$ sufficiently large for its construction and mass bounds. Put
\[
 H_j=H+\frac{\log L}{\pi}j,\qquad I_H=[0.9H,1.1H].
\]
An \emph{admitted nearest-neighbor trajectory of depth $N$}
is the exact $N$-step block integration of the nearest-neighbor
Wilson law constructed in \cite[Sections~5.1 and~7.1]{SharpO4}.
At every remaining depth $j=N,\ldots,0$, its retained density
satisfies the controlled-class bounds in that source and its
kinetic coupling lies in $[H_j/2,2H_j]$. The kinetic coupling
is the coefficient extracted from the quadratic energy of a
slowly varying retained spin field. Its initial value is the
bare coupling $\beta$ and its final value is denoted by $h$.
No uniqueness of an admitted bare coupling with given $N,h$
is assumed. The trace argument will use only the numerical
conclusions of the following imported proposition.

\begin{proposition}[Wilson trajectory input]
\label{prop:wilson-input}
For each $N\ge1$ and $h\in I_H$ there exists an admitted
nearest-neighbor trajectory of depth $N$ ending at $h$.
Every sufficiently large bare coupling occurs in at least
one such trajectory. There are fixed integers $N_0,M_0>0$,
with $M_0$ dyadic, such that every admitted trajectory with
$N\ge N_0$, terminal coupling $h\in I_H$, bare coupling
$\beta$, and spacing $a=L^{-N}$ obeys
\begin{align}
 4\log Z_\beta(M_0/a,M_0/a)
 -\log Z_\beta(2M_0/a,2M_0/a)&\le2^{-10},
 \label{eq:wilson-seed}\\
 m(\beta)/a&\le\log8,
 \label{eq:wilson-input-upper}\\
 N\log L&=\pi\beta-\tfrac12\log\beta+O_{L,H}(1).
 \label{eq:wilson-calibration}
\end{align}
In particular $a\asymp_{L,H}\sqrt\beta e^{-\pi\beta}$.
The error bound and the two preceding inequalities are uniform
over these trajectories. The periodic plane limit is unique
and has a positive full OS gap for every fixed $\beta>0$.
\end{proposition}

\begin{proof}[Source and applicability]
Admission and coverage are \cite[Proposition~7.1]{SharpO4};
calibration is its Proposition~6.2. The seed
\eqref{eq:wilson-seed} is the concluding estimate in the proof
of its Theorem~9.2: the reference depth and then the dyadic
$M_0$ are fixed before the final depth varies. The upper bound
is its Proposition~10.2. Periodic-limit uniqueness and the
all-temperature full gap are its Theorem~1.1 and Corollary~1.2.
The sphere probability convention, single counting of bonds,
even periods, and site-reflection one-step transfer coincide
with ours. In particular the cited gaps concern the full local
OS space, including rotation-invariant observables.
\end{proof}

Fix one admitted trajectory with $N\ge N_0$ and $h\in I_H$
for the following definitions and suppress its parameters.
Physical lengths $t,w$ always have $t/a,w/a$
even integers at least four. Write
\[
 \mathcal Z(t,w)=Z_\beta(t/a,w/a),\qquad
 U(w)=U_{\beta,w/a},\qquad E(w)=-a^{-1}\log\tau_{\beta,w/a}.
\]
Define
\begin{align*}
 p(t,w)&=2\log\mathcal Z(t,w)-\log\mathcal Z(2t,w),\\
 \Delta(w)&=4\log\mathcal Z(w,w)-\log\mathcal Z(2w,2w),\\
 s(t,w)&=\operatorname{Tr}U(w)^{t/a}-1.
\end{align*}
The top-eigenvalue scalar cancels in $p$, and every
area-proportional free-energy term cancels in $\Delta$.

\begin{proposition}[Uniform cylinder and partition rates]
\label{prop:cylinder-rate}
\label{prop:wilson-partition-rate}
For all admitted trajectories of Proposition~\ref{prop:wilson-input}
with $N\ge N_0$ and terminal coupling $h\in I_H$, there are
constants $C_H,w_H<\infty$, independent of the trajectory,
such that every dyadic $w\ge w_H$ satisfies
\begin{equation}
 \left|E(w)-\frac{m(\beta)}a\right|\le C_Hw^{-1/4}
 \label{eq:wilson-cylinder-rate}
\end{equation}
and
\begin{equation}
 \left|-\frac{\log\Delta(w)}w-\frac{m(\beta)}a\right|
       \le C_Hw^{-1/4}.
 \label{eq:partition-rate}
\end{equation}
\end{proposition}

We first obtain uniform convergence of a truncated cylinder
rate from exact rectangle identities. We then identify its
limit with the plane gap and remove the truncation.

\subsection{Trace tails and rectangle identities}

The transfer eigenvalues give
\[
 p(t,w)=2\log(1+s(t,w))-\log(1+s(2t,w)).
\]
If $1=u_1\ge u_2\ge\cdots\ge0$ are the eigenvalues of $U(w)$, the
normalized trace weights at time $t$ are
$\eta_i=u_i^{t/a}/(1+s(t,w))$. Thus
$e^{-p(t,w)}=\sum_i\eta_i^2\le\eta_1=(1+s(t,w))^{-1}$.
In particular $p\ge0$. Rotation interchanges the two periods; substitution
in the definitions proves
\begin{equation}
 \begin{aligned}
 \Delta(w)&=2p(w,w)+p(w,2w),\\
 p(t,2w)&=2p(t,w)-2p(w,t)+p(w,2t),\\
 0\le p(t,w)&\le2s(t,w),\\
 s(t,w)&\le e^{p(t,w)}-1.
 \end{aligned}
 \label{eq:rectangle-identities}
\end{equation}
Also,
\begin{equation}
 e^{-tE(w)}\le s(t,w)
 \le s(t_0,w)e^{-(t-t_0)E(w)}\quad(t\ge t_0).
 \label{eq:wilson-trace-advance}
\end{equation}
If $p(t,w)\le1/4$, then $s\le pe^p$ gives
\begin{equation}
 e^{-tE(w)}\le2p(t,w),\qquad
 p(2t,w)\le4p(t,w)^2.
 \label{eq:wilson-squaring}
\end{equation}
For the second inequality use
$p(2t,w)\le2s(2t,w)\le2s(t,w)^2\le2p^2e^{2p}$.

We include the dyadic argument from
\cite[Proposition~2.2]{SharpO4}. If $\Delta(r)\le2^{-10}$,
then $p(r,2r)\le\Delta(r)$ and
$p(2r,2r)\le4\Delta(r)^2$. The rectangle identity gives
\[
 p(r,4r)\le2p(r,2r)+p(2r,2r)\le3\Delta(r).
\]
Squaring yields $p(2r,4r)\le36\Delta(r)^2$, hence
\[
 \Delta(2r)\le44\Delta(r)^2\le64\Delta(r)^2.
\]
All squaring steps are justified by $3\cdot2^{-10}<1/4$.
Starting with \eqref{eq:wilson-seed}, this proves on dyadic
$w=2^kM_0$, with uniform $c,C>0$,
\begin{equation}
 \Delta(w)\le Ce^{-cw},\qquad E(w)\ge c,\qquad
 s(w/2,w)\le Ce^{-cw}\quad(w\ge2M_0).
 \label{eq:wilson-basic-tails}
\end{equation}
Indeed $e^{-wE(w)}\le2p(w,w)\le\Delta(w)$, and
$p(w/2,w)\le\Delta(w/2)$ controls the half-time trace.

\begin{lemma}[Long rectangles]
\label{lem:wilson-long-rectangles}
Uniformly over admitted trajectories with $N\ge N_0$ and
$h\in I_H$, for all sufficiently large dyadic $t$ and dyadic $w\ge2t$,
\begin{equation}
 p(t,w)\le C(w/t)e^{-ct}.
 \label{eq:wilson-long-rectangle}
\end{equation}
\end{lemma}

\begin{proof}
At $w=2t$, $p(t,2t)\le\Delta(t)$. For widths $z=t,2t$,
advance the half-circumference estimate
\eqref{eq:wilson-basic-tails} to time $w$ using
\eqref{eq:wilson-trace-advance} and $E(z)\ge c$.
This gives $p(w,t)+p(w,2t)\le Ce^{-cw}$, after decreasing $c$.
The rectangle identity therefore gives
$p(t,2w)\le2p(t,w)+Ce^{-cw}$.
Iteration over dyadic widths proves the result.
\end{proof}

\subsection{Convergence of a truncated cylinder rate}

Choose dyadic $t\asymp w^{1/2}$ and $t_0\asymp w^{1/4}$,
with fixed implicit constants. At widths $W=w,2w$,
Lemma~\ref{lem:wilson-long-rectangles} makes
$s(t_0,W)\le1$ uniformly for all large $w$. Hence
\begin{equation}
 e^{-tE(W)}\le s(t,W)\le e^{-(t-t_0)E(W)},\qquad
 \tfrac12s(t,W)\le p(t,W)\le2s(t,W).
 \label{eq:wilson-short-time-sandwich}
\end{equation}
The lower bound uses $p\ge\log(1+s)\ge s/2$ when $s\le1$.

We do not yet have a uniform upper bound for $E(W)$.
For fixed $B>0$, set
\[
 f_B(t,W)=-\frac1t\log\bigl(p(t,W)+e^{-Bt}\bigr).
\]
The preceding sandwich proves
\begin{equation}
 f_B(t,W)=\min\{E(W),B\}+O_B(w^{-1/4}).
 \label{eq:wilson-capped-rate}
\end{equation}
Its upper and lower exponential rates are
$\min\{E(W),B\}$ and
$\min\{(1-t_0/t)E(W),B\}$, up to $O(1/t)$.
Their difference is at most $Bt_0/t$, even if $E(W)>B$.

The swapped terms in the rectangle identity are $O(e^{-cw})$,
so $p(t,2w)=2p(t,w)+O(e^{-cw})$. The added $e^{-Bt}$ bounds
both logarithm arguments below, and $t\asymp\sqrt w$ makes
the additive error negligible. The factor two changes their
logarithms by at most a constant. Thus
$|f_B(t,2w)-f_B(t,w)|\le C_B/t$, and
\[
 |\min\{E(2w),B\}-\min\{E(w),B\}|\le C_Bw^{-1/4}.
\]
This is summable along dyadic doublings. The truncated rate
has a limit for each trajectory, with convergence error
$O_B(w^{-1/4})$ uniform over all trajectories.

\subsection{Identification with the plane gap}

For each fixed trajectory we first prove
\begin{equation}
 m(\beta)/a\ge\limsup_{w\to\infty}E(w),
 \label{eq:wilson-first-gap-direction}
\end{equation}
with dyadic widths. If $F$ is a bounded continuous observable
in a time slab of $r$ lattice steps, its normalized transfer
insertion $A_F$ has norm at most $\|F\|_\infty$. For $r=0$ it is
multiplication by $F$; for $r\ge1$ its kernel is pointwise dominated
in modulus by $\|F\|_\infty U(w)^r$.
For two such insertions, let $d$ be the number of lattice steps from
the last row of the first slab to the first row of the second. Then
\[
 |\operatorname{Cov}_{\omega_{\beta,w/a}}(F,G)|
 \le\|F\|_\infty\|G\|_\infty e^{-adE(w)}.
\]
The square-torus expectation of $F$ is
\[
 \frac{\operatorname{Tr}(A_F U(w)^{\,w/a-r})}
      {\operatorname{Tr}U(w)^{\,w/a}}.
\]
When $ar\le w/2$, separating the ground projection in
numerator and denominator bounds its difference from the
cylinder expectation by $2\|F\|_\infty s(w/2,w)$.
By \eqref{eq:wilson-basic-tails}, cylinders and square tori
therefore have the same local limits, namely the unique
periodic plane state.

Take a subsequence realizing the limsup of $E(w)$ and pass
the slab inequality to the plane. Applied to reflected local
autocorrelations, positivity of their spectral measures and
density in the OS space prove
\eqref{eq:wilson-first-gap-direction}. An infinite limsup would
give an infinite gap by the same argument, contradicting
\eqref{eq:wilson-input-upper}.

Take $B=10$ in the truncated-rate convergence. Its limit
is at most $\log8<10$, by
\eqref{eq:wilson-first-gap-direction} and
\eqref{eq:wilson-input-upper}. Uniform convergence makes the
truncation inactive for all sufficiently large widths,
uniformly over the trajectories. Thus
\begin{equation}
 E_\infty=\lim_{w\to\infty}E(w),\qquad
 |E(w)-E_\infty|\le C_Hw^{-1/4},\qquad
 E_\infty\le m(\beta)/a\le\log8.
 \label{eq:wilson-cylinder-limit}
\end{equation}

For the reverse inequality, fix $\beta,a$ and keep
$t\asymp\sqrt w$. Apply Lemmas~\ref{lem:wilson-detector}
and~\ref{lem:wilson-visibility} at $n=t/a$.
Since $E(w)$ is bounded, $\tau_{\beta,w/a}=e^{-aE(w)}$
and the detector overlap are bounded below by positive
constants for this fixed trajectory. Some pair on the two
rows therefore has correlation
\begin{equation}
 \omega_{\beta,w/a}(q_x^1q_y^1)
 \ge c_\beta'(w/a)^{-2}e^{-tE(w)}.
 \label{eq:wilson-visible-pair}
\end{equation}
The positive constant is independent of $w$.

Lift the spatial pair along a shortest circle arc. Its spatial
span is at most $w/2$, and its temporal span $t$ is also at most
$w/2$ for large $w$. The cylinder expectation and its plane
counterpart satisfy
\begin{equation}
 \left|\omega_{\beta,w/a}(q_x^1q_y^1)
       -\mathbb E_{\mu_\beta}(q_{\widetilde x}^1
                             q_{\widetilde y}^1)\right|
 \le Ce^{-cw},
 \label{eq:wilson-pair-to-plane}
\end{equation}
as follows. First compare the cylinder to the $w$ by $w$
torus by the slab estimate, with error at most $2s(w/2,w)$.
Compare the $w$ by $w$ and $2w$ by $w$ tori through their
common width-$w$ cylinder. Next rotate and compare the
$2w$ by $w$ and $2w$ by $2w$ tori through the width-$2w$
cylinder. The latter error is bounded by a constant times
$s(w/2,2w)=O(e^{-cw})$, using
Lemma~\ref{lem:wilson-long-rectangles}. Translations fit the
product into the required arcs at every comparison.
Repeat these two-direction doublings. Their errors sum to
$Ce^{-cw}$. For each initial $w$, the lifted observable is
fixed throughout the larger-torus limit, so periodic local
convergence proves \eqref{eq:wilson-pair-to-plane}.

Because $t\asymp\sqrt w$ and $E(w)$ is bounded, this error
is negligible relative to the right side of
\eqref{eq:wilson-visible-pair} as $w\to\infty$ at the fixed
trajectory. The plane component pair is at most
$e^{-m(\beta)t/a}$ by Proposition~\ref{prop:wilson-full-gap},
regardless of its spatial separation. Taking logarithms gives
\[
 m(\beta)/a
 \le E(w)+\frac{2\log(w/a)-\log(c_\beta'/2)}{t}.
\]
The extra term tends to zero, so $m(\beta)/a\le E_\infty$.
Together with \eqref{eq:wilson-cylinder-limit}, this proves
equality and \eqref{eq:wilson-cylinder-rate}. The fixed-trajectory
argument removed the nonuniform visibility constants; the
convergence rate was already uniform.

\begin{proof}[Completion of Proposition~\ref{prop:cylinder-rate}]
Apply \eqref{eq:wilson-trace-advance} at time $w$, with
$t_0\asymp w^{1/4}$ and widths $W=w,2w$.
The long-rectangle bound gives $s(t_0,W)\le1$, and hence
\[
 e^{-wE(W)}\le s(w,W)\le e^{-(w-t_0)E(W)}.
\]
Now $E(W)$ is uniformly bounded and $p(w,W)$ lies between
$s(w,W)/2$ and $2s(w,W)$. Therefore
\[
 -w^{-1}\log p(w,W)
 =E(W)+O_H(w^{-3/4})
 =m(\beta)/a+O_H(w^{-1/4}).
\]
The positive sum $\Delta(w)=2p(w,w)+p(w,2w)$ has the same
logarithmic rate within this error. This proves
\eqref{eq:partition-rate}.
\end{proof}
